% GENERATED FILE. Edit canonical lessons, then run npm run generate:books.
% canonical-source-hash: fnv1a64:fca44634d07fb824
% canonical-lessons: PA-C185-possessive-return, PA-C185-meri, PA-C185-mere, PA-C185-tuhadi, PA-C185-apni, PA-C185-apne, PA-C185-sada, PA-C185-sade, PA-C185-possessive-review, PA-C185-payoff

\chapter{Whose Book, Whose Home?}
\label{ch:pa-a2-possessive-agreement}
\hlchaptermodality{185}

\begin{chapteropening}
\textbf{By the end of this chapter:} \emph{I can recognise and use common my, your, one's-own, and our forms that agree with familiar nouns.}

Retrieve six exact possessive forms through listening, speaking, reading, and bounded copying after a slow agreement ramp.
\end{chapteropening}

\section[(possessive base review)]{Start with three forms you know}

\label{lesson:PA-C185-possessive-return}



\begin{quote}
Say \textbf{\pa{ਮੇਰਾ}}, \textbf{\pa{ਤੁਹਾਡਾ}}, then \textbf{\pa{ਆਪਣਾ}}.
\end{quote}

\begin{grammarlens}[title={Grammar Lens: the ending will follow the owned noun}]
These familiar forms all end in \emph{-aa}. Keep their meanings steady: my, respectful your, and one's own. The next steps will change one ending at a time to agree with a familiar noun.
\end{grammarlens}

\subsection*{Your turn}
\begin{itemize}
\raggedright
  \item \emph{Hear:} \emph{meraa}; YOU SAY: \textbf{\pa{ਮੇਰਾ}}
  \item \emph{Hear:} \emph{tuhaadaa}; YOU SAY: \textbf{\pa{ਤੁਹਾਡਾ}}
  \item \emph{Hear:} \emph{aapnaa}; YOU SAY: \textbf{\pa{ਆਪਣਾ}}
  \item \emph{Read it:} \textbf{\pa{ਮੇਰਾ}, \pa{ਤੁਹਾਡਾ}, \pa{ਆਪਣਾ}}
\end{itemize}

\begin{tcolorbox}[breakable,colback=teal!4,colframe=teal!35!black,title={Before you move on}]
Which form means my? Respectful your? One's own?

This review introduces no new written target.
\end{tcolorbox}
\section[\texorpdfstring{\pa{ਮੇਰੀ} (merii)}{ਮੇਰੀ (merii)}]{\pa{ਮੇਰੀ} — my, with a feminine noun}

\label{lesson:PA-C185-meri}



\begin{quote}
Read \textbf{\pa{ਮੇਰਾ}}, then say \textbf{\pa{ਕਿਤਾਬ}}.
\end{quote}

\begin{grammarlens}[title={Grammar Lens: change the agreement ending}]
\begin{quote}
\textbf{\pa{ਮੇਰੀ}} — \emph{merii} — my, agreeing with a feminine noun
\end{quote}

The ending changes from \emph{-aa} to \emph{-ii}. Keep the owner the same. In \textbf{\pa{ਮੇਰੀ} \pa{ਕਿਤਾਬ}}, the form agrees with \textbf{\pa{ਕਿਤਾਬ}}.
\end{grammarlens}

\subsection*{Your turn}
\begin{itemize}
\raggedright
  \item \emph{Hear:} \emph{meraa, merii}; YOU SAY: \textbf{\pa{ਮੇਰਾ}, \pa{ਮੇਰੀ}}
  \item \emph{Read it:} \textbf{\pa{ਮੇਰੀ} \pa{ਕਿਤਾਬ}}
  \item \emph{Write it:} copy \textbf{\pa{ਮੇਰੀ}} once and check the final \textbf{\pa{ੀ}}
\end{itemize}

\begin{tcolorbox}[breakable,colback=teal!4,colframe=teal!35!black,title={Before you move on}]
Read the exact form in “my book.” (\textbf{\pa{ਮੇਰੀ}}.)

Source: the canonical Punjabi \textbf{\pa{ਮੇਰਾ}} entry already names \emph{merii} as its feminine agreement form; this lesson owns its exact Gurmukhi spelling.
\end{tcolorbox}
\section[\texorpdfstring{\pa{ਮੇਰੇ} (mere)}{ਮੇਰੇ (mere)}]{\pa{ਮੇਰੇ} — my, with plural friends}

\label{lesson:PA-C185-mere}



\begin{quote}
Say \textbf{\pa{ਮੇਰਾ}}, then \textbf{\pa{ਦੋਸਤ}}.
\end{quote}

\begin{grammarlens}[title={Grammar Lens: use the \emph{-e} agreement form}]
\begin{quote}
\textbf{\pa{ਮੇਰੇ}} — \emph{mere} — my, with plural or oblique masculine nouns
\end{quote}

The owner is still “me.” The ending changes to \emph{-e}. In \textbf{\pa{ਮੇਰੇ} \pa{ਦੋਸਤ}}, the form goes with friends in the plural.
\end{grammarlens}

\subsection*{Your turn}
\begin{itemize}
\raggedright
  \item \emph{Hear:} \emph{meraa, mere}; YOU SAY: \textbf{\pa{ਮੇਰਾ}, \pa{ਮੇਰੇ}}
  \item \emph{Read it:} \textbf{\pa{ਮੇਰੇ} \pa{ਦੋਸਤ}}
  \item \emph{Write it:} copy \textbf{\pa{ਮੇਰੇ}} once and check the final \textbf{\pa{ੇ}}
\end{itemize}

\begin{tcolorbox}[breakable,colback=teal!4,colframe=teal!35!black,title={Before you move on}]
Read the exact form in “my friends.” (\textbf{\pa{ਮੇਰੇ}}.)

Source: the canonical Punjabi \textbf{\pa{ਮੇਰਾ}} entry already names \emph{mere} as its plural agreement form; this lesson owns its exact Gurmukhi spelling.
\end{tcolorbox}
\section[\texorpdfstring{\pa{ਤੁਹਾਡੀ} (tuhaadii)}{ਤੁਹਾਡੀ (tuhaadii)}]{\pa{ਤੁਹਾਡੀ} — respectful your, with a feminine noun}

\label{lesson:PA-C185-tuhadi}



\begin{quote}
Retrieve \textbf{\pa{ਤੁਹਾਡਾ}} from the respectful name question. Then say \textbf{\pa{ਯਾਤਰਾ}}.
\end{quote}

\begin{grammarlens}[title={Grammar Lens: keep respect, change agreement}]
\begin{quote}
\textbf{\pa{ਤੁਹਾਡੀ}} — \emph{tuhaadii} — your, respectful, with a feminine noun
\end{quote}

The respectful owner does not change. Only the ending becomes \emph{-ii}. Read \textbf{\pa{ਤੁਹਾਡੀ} \pa{ਯਾਤਰਾ}}: “your journey.”
\end{grammarlens}

\subsection*{Your turn}
\begin{itemize}
\raggedright
  \item \emph{Hear:} \emph{tuhaadaa, tuhaadii}; YOU SAY: \textbf{\pa{ਤੁਹਾਡਾ}, \pa{ਤੁਹਾਡੀ}}
  \item \emph{Read it:} \textbf{\pa{ਤੁਹਾਡੀ} \pa{ਯਾਤਰਾ}}
  \item \emph{Write it:} copy \textbf{\pa{ਤੁਹਾਡੀ}} once and check the final \textbf{\pa{ੀ}}
\end{itemize}

\begin{tcolorbox}[breakable,colback=teal!4,colframe=teal!35!black,title={Before you move on}]
Read the respectful possessive in “your journey.” (\textbf{\pa{ਤੁਹਾਡੀ}}.)

Source: the canonical respectful \textbf{\pa{ਤੁਹਾਡਾ}} phrase and ordinary Punjabi agreement supply the feminine form used here.
\end{tcolorbox}
\section[\texorpdfstring{\pa{ਆਪਣੀ} (aapnii)}{ਆਪਣੀ (aapnii)}]{\pa{ਆਪਣੀ} — one's own, with a feminine noun}

\label{lesson:PA-C185-apni}



\begin{quote}
Say \textbf{\pa{ਆਪਣਾ}}, then \textbf{\pa{ਕਿਤਾਬ}}.
\end{quote}

\begin{grammarlens}[title={Grammar Lens: the same \emph{-ii} agreement}]
\begin{quote}
\textbf{\pa{ਆਪਣੀ}} — \emph{aapnii} — one's own, with a feminine noun
\end{quote}

The reflexive meaning stays “one's own.” The ending follows the feminine noun. Read \textbf{\pa{ਆਪਣੀ} \pa{ਕਿਤਾਬ}}.
\end{grammarlens}

\subsection*{Your turn}
\begin{itemize}
\raggedright
  \item \emph{Hear:} \emph{aapnaa, aapnii}; YOU SAY: \textbf{\pa{ਆਪਣਾ}, \pa{ਆਪਣੀ}}
  \item \emph{Read it:} \textbf{\pa{ਆਪਣੀ} \pa{ਕਿਤਾਬ}}
  \item \emph{Write it:} copy \textbf{\pa{ਆਪਣੀ}} once and check the final \textbf{\pa{ੀ}}
\end{itemize}

\begin{tcolorbox}[breakable,colback=teal!4,colframe=teal!35!black,title={Before you move on}]
Read the exact form in “one's own book.” (\textbf{\pa{ਆਪਣੀ}}.)

Source: the canonical Punjabi \textbf{\pa{ਆਪਣਾ}} entry and regular feminine agreement.
\end{tcolorbox}
\section[\texorpdfstring{\pa{ਆਪਣੇ} (aapne)}{ਆਪਣੇ (aapne)}]{\pa{ਆਪਣੇ} — one's own, with plural friends}

\label{lesson:PA-C185-apne}



\begin{quote}
Say \textbf{\pa{ਆਪਣਾ}}, then \textbf{\pa{ਦੋਸਤ}}.
\end{quote}

\begin{grammarlens}[title={Grammar Lens: use the \emph{-e} agreement form}]
\begin{quote}
\textbf{\pa{ਆਪਣੇ}} — \emph{aapne} — one's own, with plural or oblique masculine nouns
\end{quote}

The meaning stays reflexive. The ending becomes \emph{-e}. Read \textbf{\pa{ਆਪਣੇ} \pa{ਦੋਸਤ}}.
\end{grammarlens}

\subsection*{Your turn}
\begin{itemize}
\raggedright
  \item \emph{Hear:} \emph{aapnaa, aapne}; YOU SAY: \textbf{\pa{ਆਪਣਾ}, \pa{ਆਪਣੇ}}
  \item \emph{Read it:} \textbf{\pa{ਆਪਣੇ} \pa{ਦੋਸਤ}}
  \item \emph{Write it:} copy \textbf{\pa{ਆਪਣੇ}} once and check the final \textbf{\pa{ੇ}}
\end{itemize}

\begin{tcolorbox}[breakable,colback=teal!4,colframe=teal!35!black,title={Before you move on}]
Read the exact form in “one's own friends.” (\textbf{\pa{ਆਪਣੇ}}.)

Source: the canonical Punjabi \textbf{\pa{ਆਪਣਾ}} entry and regular plural agreement.
\end{tcolorbox}
\section[\texorpdfstring{\pa{ਸਾਡਾ} (saadaa)}{ਸਾਡਾ (saadaa)}]{\pa{ਸਾਡਾ} — our, the missing base}

\label{lesson:PA-C185-sada}



\begin{quote}
Say \textbf{\pa{ਅਸੀਂ}}, then \textbf{\pa{ਘਰ}}.
\end{quote}

\begin{grammarlens}[title={Grammar Lens: connect we to our}]
\begin{quote}
\textbf{\pa{ਸਾਡਾ}} — \emph{saadaa} — our, with a masculine singular noun
\end{quote}

This base was not printed in earlier lessons, so pause here before changing its ending. Read \textbf{\pa{ਸਾਡਾ} \pa{ਘਰ}}: “our home.”
\end{grammarlens}

\subsection*{Your turn}
\begin{itemize}
\raggedright
  \item \emph{Hear:} \emph{asii, saadaa}; YOU SAY: \textbf{\pa{ਅਸੀਂ}, \pa{ਸਾਡਾ}}
  \item \emph{Read it:} \textbf{\pa{ਸਾਡਾ} \pa{ਘਰ}}
  \item \emph{Write it:} copy \textbf{\pa{ਸਾਡਾ}} once
\end{itemize}

\begin{tcolorbox}[breakable,colback=teal!4,colframe=teal!35!black,title={Before you move on}]
Read the new base meaning “our.” (\textbf{\pa{ਸਾਡਾ}}.)

Source: standard Punjabi \textbf{\pa{ਸਾਡਾ}}, the masculine singular possessive paired with \textbf{\pa{ਅਸੀਂ}}.
\end{tcolorbox}
\section[\texorpdfstring{\pa{ਸਾਡੇ} (saade)}{ਸਾਡੇ (saade)}]{\pa{ਸਾਡੇ} — our, with plural homes}

\label{lesson:PA-C185-sade}



\begin{quote}
Read the new base \textbf{\pa{ਸਾਡਾ}}, then say \textbf{\pa{ਘਰ}}.
\end{quote}

\begin{grammarlens}[title={Grammar Lens: change only the ending}]
\begin{quote}
\textbf{\pa{ਸਾਡੇ}} — \emph{saade} — our, with plural or oblique masculine nouns
\end{quote}

The owner remains “we.” The ending changes from \emph{-aa} to \emph{-e}. Read \textbf{\pa{ਸਾਡੇ} \pa{ਘਰ}}: “our homes.”
\end{grammarlens}

\subsection*{Your turn}
\begin{itemize}
\raggedright
  \item \emph{Hear:} \emph{saadaa, saade}; YOU SAY: \textbf{\pa{ਸਾਡਾ}, \pa{ਸਾਡੇ}}
  \item \emph{Read it:} \textbf{\pa{ਸਾਡੇ} \pa{ਘਰ}}
  \item \emph{Write it:} copy \textbf{\pa{ਸਾਡੇ}} once and check the final \textbf{\pa{ੇ}}
\end{itemize}

\begin{tcolorbox}[breakable,colback=teal!4,colframe=teal!35!black,title={Before you move on}]
Read the exact form in “our homes.” (\textbf{\pa{ਸਾਡੇ}}.)

Source: standard Punjabi plural agreement built from the newly owned base \textbf{\pa{ਸਾਡਾ}}.
\end{tcolorbox}
\section[\texorpdfstring{\pa{ਮੇਰੀ}, \pa{ਮੇਰੇ}, \pa{ਤੁਹਾਡੀ} … (merii, mere …)}{ਮੇਰੀ, ਮੇਰੇ, ਤੁਹਾਡੀ … (merii, mere …)}]{Seven forms in six small phrases}

\label{lesson:PA-C185-possessive-review}



\begin{quote}
Say the four \emph{-ii} and \emph{-e} contrasts from memory, then the new “our” base and its \emph{-e} form.
\end{quote}

\subsection*{Your turn: keep each form beside its noun}
Read slowly:

\begin{itemize}
\raggedright
  \item \textbf{\pa{ਮੇਰੀ} \pa{ਕਿਤਾਬ}} — my book
  \item \textbf{\pa{ਮੇਰੇ} \pa{ਦੋਸਤ}} — my friends
  \item \textbf{\pa{ਤੁਹਾਡੀ} \pa{ਯਾਤਰਾ}} — your journey
  \item \textbf{\pa{ਆਪਣੀ} \pa{ਕਿਤਾਬ}} — one's own book
  \item \textbf{\pa{ਆਪਣੇ} \pa{ਦੋਸਤ}} — one's own friends
  \item \textbf{\pa{ਸਾਡਾ} \pa{ਘਰ}, \pa{ਸਾਡੇ} \pa{ਘਰ}} — our home, our homes
\end{itemize}

\emph{Hear:} the six phrases in mixed order; YOU SAY: each matching owner and noun

\begin{tcolorbox}[breakable,colback=teal!4,colframe=teal!35!black,title={Before you move on}]
\emph{Write it:} cover the models and write \textbf{\pa{ਮੇਰੀ}, \pa{ਮੇਰੇ}, \pa{ਤੁਹਾਡੀ}, \pa{ਆਪਣੀ}, \pa{ਆਪਣੇ}, \pa{ਸਾਡਾ}, \pa{ਸਾਡੇ}} once

Check the \emph{-ii} forms first, then the \emph{-e} forms and the one new \emph{-aa} base.

This review adds no new form and makes no mock-readiness claim.
\end{tcolorbox}
\section[\texorpdfstring{\pa{ਮੇਰੀ} \pa{ਕਿਤਾਬ} … (merii kitaab …)}{ਮੇਰੀ ਕਿਤਾਬ … (merii kitaab …)}]{A gentle possessive-agreement payoff}

\label{lesson:PA-C185-payoff}



\begin{quote}
Without looking, say one \emph{-ii} form, one \emph{-e} form, and the new “our” base.
\end{quote}

\subsection*{Your turn: four skills, bounded to agreement}
\begin{itemize}
\raggedright
  \item \emph{Hear:} \emph{merii kitaab, tuhaadii yaatraa, aapne dost, saade ghar}; YOU SAY: \textbf{\pa{ਮੇਰੀ} \pa{ਕਿਤਾਬ}, \pa{ਤੁਹਾਡੀ} \pa{ਯਾਤਰਾ}, \pa{ਆਪਣੇ} \pa{ਦੋਸਤ}, \pa{ਸਾਡੇ} \pa{ਘਰ}}
  \item \emph{Read it:} \textbf{\pa{ਮੇਰੇ} \pa{ਦੋਸਤ।} \pa{ਆਪਣੀ} \pa{ਕਿਤਾਬ।} \pa{ਸਾਡਾ} \pa{ਘਰ।}}
  \item \emph{Write it:} cover the model and write \textbf{\pa{ਮੇਰੀ}, \pa{ਮੇਰੇ}, \pa{ਤੁਹਾਡੀ}, \pa{ਆਪਣੀ}, \pa{ਆਪਣੇ}, \pa{ਸਾਡੇ}}
\end{itemize}

Now point to the three forms ending in \textbf{\pa{ੀ}}. Then point to the three measured forms ending in \textbf{\pa{ੇ}}.

\begin{tcolorbox}[breakable,colback=teal!4,colframe=teal!35!black,title={Before you move on}]
Retrieve the six measured forms once more, then say the supporting base \textbf{\pa{ਸਾਡਾ}}.

This is the eighth prerequisite tranche exposed by Punjabi A2 mock drafting. It does not unblock Punjabi A2 mock 1 by itself. It is not an A2 mock, not A2 readiness or attainment, and not external or human validation.

Source: the canonical Punjabi possessive bases and ordinary agreement with the familiar book, friend, journey, and home nouns.
\end{tcolorbox}
